\documentclass[11pt]{article}
\usepackage[T1]{fontenc}
\usepackage{lmodern}
\usepackage{amsmath,amssymb,amsthm,mathtools}
\usepackage{graphicx}
\usepackage[margin=1in]{geometry}
\usepackage[colorlinks=true,linkcolor=blue,citecolor=blue,urlcolor=blue]{hyperref}
\usepackage{microtype}

\newtheorem{theorem}{Theorem}
\newtheorem{lemma}[theorem]{Lemma}
\newtheorem{remark}[theorem]{Remark}
\newcommand{\N}{\mathbb N}
\newcommand{\C}{\mathbb C}
\newcommand{\supp}{\operatorname{supp}}

\title{Rapid decay survives exponential distortion\\
\large A nondegenerate amalgam answering the final open consideration in
arXiv:2509.08754}
\author{Automated mathematics discovery run \texttt{fa\_banach\_001}\\
Lane 5, candidate for human review}
\date{19 July 2026}

\begin{document}
\maketitle

\begin{center}
\fbox{\parbox{0.92\textwidth}{\textbf{Status: candidate full solution, likely valid.}
We construct a nondegenerate amalgamated free product whose universal length
has rapid decay although one vertex length is exponentially distorted.  The
construction uses countable groups with weighted lengths and does not address
the neighboring finitely generated word-length questions.}}
\end{center}

\section{Source question}

Kunnawalkam Elayavalli, Patchell and Teryoshin develop subexponential decay
(SD), prove permanence theorems for amalgamated free products, and quantify
the effect of distortion \cite{KPT}.  They close their introduction with three
open directions.  The last asks whether there are systematic or specific
instances in which RD or SD survives high distortion:

\begin{center}
\includegraphics[width=0.98\textwidth]{figures/open_question_crop.png}
\end{center}

We answer this final question affirmatively with a fully explicit example.
The setting is exactly the source's general framework of countable groups
equipped with integer-valued length functions.

\section{Construction and result}

Let
\[
                 A=\bigoplus_{n\ge0}C_2
\]
and denote its standard basis by $e_0,e_1,\ldots$.  Every $a\in A$ has a
unique finite support $S(a)\subset\N\cup\{0\}$.  Define two weighted lengths
on $A$ by
\begin{equation}\label{eq:lengths}
 \ell_0(a)=\sum_{n\in S(a)}2^n,
 \qquad
 \ell_1(a)=\sum_{n\in S(a)}2^{2^n}.
\end{equation}
For $u\in C_k$, put $\delta_k(u)=0$ if $u=1$ and $\delta_k(u)=1$ otherwise.
Set
\begin{align}\label{eq:vertices}
 G&=A\times C_2,& L_G(a,u)&=\ell_1(a)+\delta_2(u),\notag\\
 H&=A\times C_3,& L_H(a,v)&=\ell_0(a)+\delta_3(v).
\end{align}
We identify $A$ with $A\times\{1\}$ in both factors and form
$\Gamma=G*_AH$.  Both inclusions of $A$ are proper.

Let $L^U$ be the universal length on $\Gamma$ defined in \cite{KPT}: in a
factorization by elements of $G\cup H$, each letter is charged its vertex
length, while a letter in $A$ may be charged either restriction.

\begin{theorem}[RD beyond the polynomial-distortion regime]\label{thm:main}
There is a canonical isomorphism
\begin{equation}\label{eq:group-id}
                 \Gamma\cong A\times(C_2*C_3).
\end{equation}
If $|q|$ denotes the standard free-product word length on $C_2*C_3$, then
\begin{equation}\label{eq:universal}
                 L^U(a,q)=\ell_0(a)+|q|.
\end{equation}
The pair $(\Gamma,L^U)$ has rapid decay.

On the other hand, define the distortion function of the first vertex by
\[
 D(r)=\max\{L_G(g):g\in G,\ L^U(g)\le r\}.
\]
For every integer $r\ge1$,
\begin{equation}\label{eq:distortion}
                   2^{r/2}<D(r)\le2^{r+1}+1.
\end{equation}
Thus $(L^U|_G,L_G)$ is exponentially distorted.  In particular, the
polynomial-distortion hypothesis in the source's RD permanence theorem fails.
\end{theorem}

\section{Proof intuition}

The Boolean group is assigned two radically different proper geometries.
Under the light weights $2^n$, binary expansion gives exactly one group
element of every nonnegative integer length, so balls have linear size.  The
heavy weights $2^{2^n}$ are exponentially larger but only make balls smaller.

The amalgam is transparent because $A$ is central in both vertices: it factors
as $A$ times the virtually free group $C_2*C_3$.  In the universal length, an
$A$ coordinate can always be moved to the lighter vertex, so the light length
wins exactly.  This gives RD by direct-product permanence.  Seen from the
heavy vertex, however, the binary element of light length $r$ has heavy length
at least $2^{r/2}$.

\section{Proof of Theorem~\ref{thm:main}}

\subsection{Lengths and the amalgam}

The two functions in \eqref{eq:lengths} are weighted word lengths.  Indeed,
addition in $A$ takes symmetric differences of supports, so cancellation can
only decrease the sum of positive weights.  Hence $L_G$ and $L_H$ are length
functions.

The group $A$ is central in both $G$ and $H$.  The universal property of the
amalgam gives mutually inverse homomorphisms between $G*_AH$ and
$A\times(C_2*C_3)$, proving \eqref{eq:group-id}.

\subsection{Exact computation of the universal length}

The upper bound in \eqref{eq:universal} is immediate: represent $a$ as an
element of the $H$ copy, at cost $\ell_0(a)$, and spell the reduced word $q$,
at cost $|q|$.

For the reverse bound, consider any factorization of $(a,q)$ by letters
$x_i\in G\cup H$.  Write the $A$ coordinate of $x_i$ as $a_i$ and its finite
cyclic coordinate as $c_i$.  Every allowed charge for $x_i$ is at least
\[
                       \ell_0(a_i)+\delta(c_i),
\]
because $\ell_1\ge\ell_0$.  Therefore the total charge is at least
\[
 \sum_i\ell_0(a_i)+\sum_i\delta(c_i)
 \ge \ell_0(a)+|q|.
\]
The last inequality uses the triangle inequality first in $A$ and then for
the word length in $C_2*C_3$.  Taking the infimum over all factorizations proves
\eqref{eq:universal}.  Notice that this also rules out savings obtained by
bundling an $A$ coordinate and a cyclic coordinate into the same factor letter.

\subsection{Rapid decay}

The map
\[
             a\longmapsto\ell_0(a)=\sum_{n\in S(a)}2^n
\]
is the binary-expansion bijection from $A$ to the nonnegative integers.  Hence
\begin{equation}\label{eq:ball}
                   |B_{\ell_0}(R)|=R+1\qquad(R\in\N\cup\{0\}).
\end{equation}
If $\varphi\in\C[A]$ is supported in this ball, then
\[
 \|\varphi\|_{C_r^*(A)}\le\|\varphi\|_1
 \le\sqrt{R+1}\,\|\varphi\|_2.
\]
Thus $(A,\ell_0)$ has RD.  Since $\ell_1\ge\ell_0$, its balls are smaller and
the same polynomial estimate proves RD for $(A,\ell_1)$.  The finite cyclic
coordinates in \eqref{eq:vertices} only multiply ball cardinalities by fixed
constants, so both vertex groups have RD.

The source's free-product Lemma 2.12 shows that $C_2*C_3$, with $|\cdot|$, has
RD: finite groups have constant decay functions, and the lemma multiplies the
maximum input function by the polynomial $(2R+1)^{7/2}$.  Its direct-product
Lemma 2.11 multiplies decay functions.  Applying it to $(A,\ell_0)$ and
$(C_2*C_3,|\cdot|)$, and using \eqref{eq:universal}, proves RD for
$(\Gamma,L^U)$.

\subsection{Exponential distortion}

For each integer $r\ge1$, let $a_r\in A$ be the element whose support is the
binary expansion of $r$.  Then $\ell_0(a_r)=r$.  If $2^N$ is its leading
binary digit, then $2^N>r/2$, so
\begin{equation}\label{eq:lower-dist}
 L_G(a_r)=\ell_1(a_r)\ge2^{2^N}>2^{r/2}.
\end{equation}
This proves the lower bound in \eqref{eq:distortion}.

Conversely, suppose $g=(a,u)\in G$ and $L^U(g)\le r$.  If $2^N$ is the
largest light weight in $S(a)$, then $2^N\le\ell_0(a)\le r$.  The
superincreasing heavy weights satisfy
\[
 \ell_1(a)\le\sum_{n=0}^N2^{2^n}
 \le2\,2^{2^N}\le2^{r+1}.
\]
Adding $\delta_2(u)\le1$ proves the upper bound in
\eqref{eq:distortion} and completes the proof.

\section{Scope and novelty check}

This is a full affirmative answer to the final request for a specific
high-distortion amalgam preserving RD/SD.  It works in the paper's stated
countable-group framework and is nondegenerate, since $A$ has indices two and
three in the vertex groups.  The example deliberately exploits weighted
lengths.  It does not answer the separate preceding questions about finitely
generated or finitely presented examples with word length.

The run's cheap indexes were searched by arXiv identifier, exact title, and
the core rapid-decay, distortion and amalgam phrases.  A bounded primary arXiv
search found the source paper but no later resolution or occurrence of this
weighted Boolean construction.  This is not an exhaustive novelty
certification.

\begin{thebibliography}{9}
\bibitem{KPT}
S. Kunnawalkam Elayavalli, G. Patchell and L. Teryoshin,
\emph{Some remarks on decay in countable groups and amalgamated free products},
arXiv:2509.08754 (2025).
\end{thebibliography}

\end{document}
